
RLHF-trained language models achieve high aggregate benchmark accuracy yet exhibit systematic calibration inversion---confidently predicting wrong answers---on specific task clusters. This paper traces this failure pattern to reward signal conflation during training, investigating a mechanism where annotator behavior propagates through reward models into model representations.

\subsubsection{The Calibration Paradox}

Standard evaluation reports aggregate accuracy, masking task-level failure patterns. When benchmark tasks are clustered by calibration scores---measuring the gap between model confidence and correctness---non-random structure emerges. Tasks where models show P(wrong) > P(correct) + 0.1 cluster with silhouette score 0.6016, exceeding the 0.3 threshold indicating meaningful clustering. This systematic pattern suggests underlying causes beyond random error.

\subsubsection{From Symptoms to Mechanisms}

Surface-level observations note that RLHF models show overconfidence on certain tasks. The deeper problem is that existing evaluation provides no mechanism explanation---we observe \textit{what} fails but not \textit{why}. Existing work lacks methods connecting calibration patterns to training dynamics.

This work tests the hypothesis that RLHF's reward modeling conflates distinct task dimensions. Specifically, if annotators rate both ''correct output'' and ''output requiring user-state modeling'' with similar confidence, the training signal fails to distinguish task types. Models then optimize for this conflated signal, potentially developing miscalibrated confidence on tasks requiring nuanced adaptation.

\subsubsection{Key Insight: The Conflation Chain}

The investigation examines a four-step mechanism:

\begin{enumerate}
\item \textbf{Annotator Conflation:} Human annotators rate correctness and user-modeling tasks identically (rate\_diff = 0.001), providing no distinguishing signal during preference data collection.
\end{enumerate}

\begin{enumerate}
\item \textbf{Reward Model Inheritance:} The reward model learns this conflated signal, showing similar confidence across task types (mean\_diff = 0.018 across three models).
\end{enumerate}

\begin{enumerate}
\item \textbf{Representation Conflation:} Model hidden states fail to separate task types internally (separation = 0.024), indicating the conflation propagates to learned representations.
\end{enumerate}

\begin{enumerate}
\item \textbf{Calibration Consequences:} These models exhibit systematic calibration inversion on specific task clusters.
\end{enumerate}

Steps 1-3 are verified empirically. Step 4's connection to bidirectional task features remains unverified with current keyword-based detection methods (2.3\% feature prevalence), motivating semantic detection as future work.

\subsubsection{Contributions}

This work makes three contributions:

\begin{itemize}
\item \textbf{Calibration clustering as diagnostic tool:} Clustering benchmark tasks by calibration inversion scores reveals systematic RLHF behavioral patterns (silhouette = 0.6016).
\end{itemize}

\begin{itemize}
\item \textbf{Empirical mechanism verification:} Quantitative evidence that annotator conflation (rate\_diff = 0.001), reward conflation (mean\_diff = 0.018), and representation conflation (separation = 0.024) form a coherent mechanism chain.
\end{itemize}

\begin{itemize}
\item \textbf{Negative result:} Keyword-based bidirectional feature detection is insufficient (2.3\% prevalence, r = -0.027), establishing methodological boundaries for future work.
\end{itemize}
